
\documentclass{article}
\usepackage{colortbl}
\usepackage{makecell}
\usepackage{multirow}
\usepackage{supertabular}

\begin{document}

\newcounter{utterance}

\centering \large Interaction Transcript for game `matchit\_ascii', experiment `similar\_grid\_2', episode 1 with deepseek{-}v4{-}pro.
\vspace{24pt}

{ \footnotesize  \setcounter{utterance}{1}
\setlength{\tabcolsep}{0pt}
\begin{supertabular}{c@{$\;$}|p{.15\linewidth}@{}p{.15\linewidth}p{.15\linewidth}p{.15\linewidth}p{.15\linewidth}p{.15\linewidth}}
    \# & $\;$A & \multicolumn{4}{c}{Game Master} & $\;\:$B\\
    \hline

    \theutterance \stepcounter{utterance}  
    & & \multicolumn{4}{p{0.6\linewidth}}{
        \cellcolor[rgb]{0.9,0.9,0.9}{
            \makecell[{{p{\linewidth}}}]{
                \texttt{\tiny{[P1$\langle$GM]}}
                \texttt{Bierzesz udział w grze polegającej na wspólnym zgadywaniu.} \\
\texttt{Celem jest ustalenie, czy ta siatka i inny obrazek, który widzę tylko ja, są takie same.} \\
\texttt{Proszę najpierw opisać swoją siatkę. Następnie ja przedstawię swój opis i będziemy mogli zadawać sobie nawzajem pytania dotyczące siatek, aby ustalić, czy są takie same. Teraz zacznij swój krótki opis siatki od "OPIS:", a następnie podaj opis. Nie dodawaj nic więcej. To jest twoja siatka:} \\
\texttt{▢ ▢ X ▢ ▢} \\
\texttt{▢ ▢ X ▢ ▢} \\
\texttt{▢ ▢ X ▢ ▢} \\
\texttt{▢ ▢ X ▢ ▢} \\
\texttt{X X X X X} \\
            }
        }
    }
    & \\ \\

    \theutterance \stepcounter{utterance}  
    & \multicolumn{4}{p{0.6\linewidth}}{
        \cellcolor[rgb]{0.9,0.9,0.9}{
            \makecell[{{p{\linewidth}}}]{
                \texttt{\tiny{[P1$\rangle$GM]}}
                \texttt{OPIS: Siatka 5x5. W trzeciej kolumnie od lewej, w pierwszych czterech wierszach od góry znajduje się symbol X. W piątym (dolnym) wierszu symbol X znajduje się we wszystkich pięciu kolumnach. Wszystkie pozostałe pola zawierają symbol ▢.} \\
            }
        }
    }
    & & \\ \\

    \theutterance \stepcounter{utterance}  
    & & & \multicolumn{2}{p{0.3\linewidth}}{
        \cellcolor[rgb]{0.9,0.9,0.9}{
            \makecell[{{p{\linewidth}}}]{
                \texttt{\tiny{[GM$|$GM]}}
                \texttt{continue} \\
            }
        }
    }
    & & \\ \\

    \theutterance \stepcounter{utterance}  
    & & \multicolumn{4}{p{0.6\linewidth}}{
        \cellcolor[rgb]{0.9,0.9,0.9}{
            \makecell[{{p{\linewidth}}}]{
                \texttt{\tiny{[GM$\rangle$P2]}}
                \texttt{Bierzesz udział w grze polegającej na wspólnym zgadywaniu.} \\
\texttt{Celem jest ustalenie, czy ta siatka i inny obrazek, który widzę tylko ja, są takie same.} \\
\texttt{Proszę najpierw opisać swoją siatkę. Następnie ja przedstawię swój opis i będziemy mogli zadawać sobie nawzajem pytania dotyczące siatek, aby ustalić, czy są takie same. Teraz zacznij swój krótki opis siatki od "OPIS:", a następnie podaj opis. Nie dodawaj nic więcej. To jest twoja siatka:} \\
\texttt{▢ ▢ X ▢ ▢} \\
\texttt{▢ ▢ X ▢ ▢} \\
\texttt{▢ ▢ ▢ ▢ ▢} \\
\texttt{▢ ▢ X X ▢} \\
\texttt{X X X X X} \\
            }
        }
    }
    & \\ \\

    \theutterance \stepcounter{utterance}  
    & & & \multicolumn{4}{p{0.6\linewidth}}{
        \cellcolor[rgb]{0.9,0.9,0.9}{
            \makecell[{{p{\linewidth}}}]{
                \texttt{\tiny{[GM$\langle$P2]}}
                \texttt{OPIS: Siatka ma 5 wierszy i 5 kolumn. W pierwszym wierszu X jest na trzeciej pozycji. W drugim wierszu X jest na trzeciej pozycji. Trzeci wiersz jest pusty. W czwartym wierszu X są na trzeciej i czwartej pozycji. W piątym wierszu wszystkie pięć pozycji to X.} \\
            }
        }
    }
     \\ \\

    \theutterance \stepcounter{utterance}  
    & & & \multicolumn{2}{p{0.3\linewidth}}{
        \cellcolor[rgb]{0.9,0.9,0.9}{
            \makecell[{{p{\linewidth}}}]{
                \texttt{\tiny{[GM$|$GM]}}
                \texttt{continue} \\
            }
        }
    }
    & & \\ \\

    \theutterance \stepcounter{utterance}  
    & & \multicolumn{4}{p{0.6\linewidth}}{
        \cellcolor[rgb]{0.9,0.9,0.9}{
            \makecell[{{p{\linewidth}}}]{
                \texttt{\tiny{[GM$\rangle$P2]}}
                \texttt{To jest mój OPISOPIS: Siatka 5x5. W trzeciej kolumnie od lewej, w pierwszych czterech wierszach od góry znajduje się symbol X. W piątym (dolnym) wierszu symbol X znajduje się we wszystkich pięciu kolumnach. Wszystkie pozostałe pola zawierają symbol ▢.} \\
\texttt{Teraz zadaj pytanie, aby dowiedzieć się o nowych aspektach mojej siatki, które mogą różnić się od twojej siatki. Zacznij od "PYTANIE:" i nie dodawaj nic więcej.} \\
            }
        }
    }
    & \\ \\

    \theutterance \stepcounter{utterance}  
    & & & \multicolumn{4}{p{0.6\linewidth}}{
        \cellcolor[rgb]{0.9,0.9,0.9}{
            \makecell[{{p{\linewidth}}}]{
                \texttt{\tiny{[GM$\langle$P2]}}
                \texttt{PYTANIE: Czy w czwartym wierszu od góry, oprócz X w trzeciej kolumnie, znajduje się także X w czwartej kolumnie?} \\
            }
        }
    }
     \\ \\

    \theutterance \stepcounter{utterance}  
    & & & \multicolumn{2}{p{0.3\linewidth}}{
        \cellcolor[rgb]{0.9,0.9,0.9}{
            \makecell[{{p{\linewidth}}}]{
                \texttt{\tiny{[GM$|$GM]}}
                \texttt{continue} \\
            }
        }
    }
    & & \\ \\

    \theutterance \stepcounter{utterance}  
    & & \multicolumn{4}{p{0.6\linewidth}}{
        \cellcolor[rgb]{0.9,0.9,0.9}{
            \makecell[{{p{\linewidth}}}]{
                \texttt{\tiny{[P1$\langle$GM]}}
                \texttt{To jest mój OPISOPIS: Siatka ma 5 wierszy i 5 kolumn. W pierwszym wierszu X jest na trzeciej pozycji. W drugim wierszu X jest na trzeciej pozycji. Trzeci wiersz jest pusty. W czwartym wierszu X są na trzeciej i czwartej pozycji. W piątym wierszu wszystkie pięć pozycji to X.} \\
\texttt{PYTANIE: Czy w czwartym wierszu od góry, oprócz X w trzeciej kolumnie, znajduje się także X w czwartej kolumnie?Rozpocznij odpowiedź od "ODPOWIEDŹ:" i nie dodawaj nic więcej.} \\
            }
        }
    }
    & \\ \\

    \theutterance \stepcounter{utterance}  
    & \multicolumn{4}{p{0.6\linewidth}}{
        \cellcolor[rgb]{0.9,0.9,0.9}{
            \makecell[{{p{\linewidth}}}]{
                \texttt{\tiny{[P1$\rangle$GM]}}
                \texttt{ODPOWIEDŹ: Nie, w czwartym wierszu od góry X znajduje się tylko w trzeciej kolumnie.} \\
            }
        }
    }
    & & \\ \\

    \theutterance \stepcounter{utterance}  
    & & & \multicolumn{2}{p{0.3\linewidth}}{
        \cellcolor[rgb]{0.9,0.9,0.9}{
            \makecell[{{p{\linewidth}}}]{
                \texttt{\tiny{[GM$|$GM]}}
                \texttt{continue} \\
            }
        }
    }
    & & \\ \\

    \theutterance \stepcounter{utterance}  
    & & \multicolumn{4}{p{0.6\linewidth}}{
        \cellcolor[rgb]{0.9,0.9,0.9}{
            \makecell[{{p{\linewidth}}}]{
                \texttt{\tiny{[P1$\langle$GM]}}
                \texttt{Teraz zadaj pytanie, aby dowiedzieć się o nowych aspektach mojej siatki, które mogą różnić się od twojej siatki. Zacznij od "PYTANIE:" i nie dodawaj nic więcej.} \\
            }
        }
    }
    & \\ \\

    \theutterance \stepcounter{utterance}  
    & \multicolumn{4}{p{0.6\linewidth}}{
        \cellcolor[rgb]{0.9,0.9,0.9}{
            \makecell[{{p{\linewidth}}}]{
                \texttt{\tiny{[P1$\rangle$GM]}}
                \texttt{PYTANIE: Czy w trzecim wierszu od góry znajduje się jakikolwiek symbol X?} \\
            }
        }
    }
    & & \\ \\

    \theutterance \stepcounter{utterance}  
    & & & \multicolumn{2}{p{0.3\linewidth}}{
        \cellcolor[rgb]{0.9,0.9,0.9}{
            \makecell[{{p{\linewidth}}}]{
                \texttt{\tiny{[GM$|$GM]}}
                \texttt{continue} \\
            }
        }
    }
    & & \\ \\

    \theutterance \stepcounter{utterance}  
    & & \multicolumn{4}{p{0.6\linewidth}}{
        \cellcolor[rgb]{0.9,0.9,0.9}{
            \makecell[{{p{\linewidth}}}]{
                \texttt{\tiny{[GM$\rangle$P2]}}
                \texttt{ODPOWIEDŹ: Nie, w czwartym wierszu od góry X znajduje się tylko w trzeciej kolumnie.} \\
\texttt{PYTANIE: Czy w trzecim wierszu od góry znajduje się jakikolwiek symbol X?Rozpocznij odpowiedź od "ODPOWIEDŹ:" i nie dodawaj nic więcej.} \\
            }
        }
    }
    & \\ \\

    \theutterance \stepcounter{utterance}  
    & & & \multicolumn{4}{p{0.6\linewidth}}{
        \cellcolor[rgb]{0.9,0.9,0.9}{
            \makecell[{{p{\linewidth}}}]{
                \texttt{\tiny{[GM$\langle$P2]}}
                \texttt{ODPOWIEDŹ: Nie, w trzecim wierszu od góry nie ma żadnego X {-} wszystkie pola są puste (▢).} \\
            }
        }
    }
     \\ \\

    \theutterance \stepcounter{utterance}  
    & & & \multicolumn{2}{p{0.3\linewidth}}{
        \cellcolor[rgb]{0.9,0.9,0.9}{
            \makecell[{{p{\linewidth}}}]{
                \texttt{\tiny{[GM$|$GM]}}
                \texttt{continue} \\
            }
        }
    }
    & & \\ \\

    \theutterance \stepcounter{utterance}  
    & & \multicolumn{4}{p{0.6\linewidth}}{
        \cellcolor[rgb]{0.9,0.9,0.9}{
            \makecell[{{p{\linewidth}}}]{
                \texttt{\tiny{[GM$\rangle$P2]}}
                \texttt{Teraz zadaj pytanie, aby dowiedzieć się o nowych aspektach mojej siatki, które mogą różnić się od twojej siatki. Zacznij od "PYTANIE:" i nie dodawaj nic więcej.} \\
            }
        }
    }
    & \\ \\

    \theutterance \stepcounter{utterance}  
    & & & \multicolumn{4}{p{0.6\linewidth}}{
        \cellcolor[rgb]{0.9,0.9,0.9}{
            \makecell[{{p{\linewidth}}}]{
                \texttt{\tiny{[GM$\langle$P2]}}
                \texttt{PYTANIE: Czy w twojej siatce w piątym (dolnym) wierszu X znajduje się we wszystkich pięciu kolumnach?} \\
            }
        }
    }
     \\ \\

    \theutterance \stepcounter{utterance}  
    & & & \multicolumn{2}{p{0.3\linewidth}}{
        \cellcolor[rgb]{0.9,0.9,0.9}{
            \makecell[{{p{\linewidth}}}]{
                \texttt{\tiny{[GM$|$GM]}}
                \texttt{continue} \\
            }
        }
    }
    & & \\ \\

    \theutterance \stepcounter{utterance}  
    & & \multicolumn{4}{p{0.6\linewidth}}{
        \cellcolor[rgb]{0.9,0.9,0.9}{
            \makecell[{{p{\linewidth}}}]{
                \texttt{\tiny{[P1$\langle$GM]}}
                \texttt{ODPOWIEDŹ: Nie, w trzecim wierszu od góry nie ma żadnego X {-} wszystkie pola są puste (▢).} \\
\texttt{PYTANIE: Czy w twojej siatce w piątym (dolnym) wierszu X znajduje się we wszystkich pięciu kolumnach?Rozpocznij odpowiedź od "ODPOWIEDŹ:" i nie dodawaj nic więcej.} \\
            }
        }
    }
    & \\ \\

    \theutterance \stepcounter{utterance}  
    & \multicolumn{4}{p{0.6\linewidth}}{
        \cellcolor[rgb]{0.9,0.9,0.9}{
            \makecell[{{p{\linewidth}}}]{
                \texttt{\tiny{[P1$\rangle$GM]}}
                \texttt{ODPOWIEDŹ: Tak, w piątym (dolnym) wierszu X znajduje się we wszystkich pięciu kolumnach.} \\
            }
        }
    }
    & & \\ \\

    \theutterance \stepcounter{utterance}  
    & & & \multicolumn{2}{p{0.3\linewidth}}{
        \cellcolor[rgb]{0.9,0.9,0.9}{
            \makecell[{{p{\linewidth}}}]{
                \texttt{\tiny{[GM$|$GM]}}
                \texttt{continue} \\
            }
        }
    }
    & & \\ \\

    \theutterance \stepcounter{utterance}  
    & & \multicolumn{4}{p{0.6\linewidth}}{
        \cellcolor[rgb]{0.9,0.9,0.9}{
            \makecell[{{p{\linewidth}}}]{
                \texttt{\tiny{[P1$\langle$GM]}}
                \texttt{Teraz zadaj pytanie, aby dowiedzieć się o nowych aspektach mojej siatki, które mogą różnić się od twojej siatki. Zacznij od "PYTANIE:" i nie dodawaj nic więcej.} \\
            }
        }
    }
    & \\ \\

    \theutterance \stepcounter{utterance}  
    & \multicolumn{4}{p{0.6\linewidth}}{
        \cellcolor[rgb]{0.9,0.9,0.9}{
            \makecell[{{p{\linewidth}}}]{
                \texttt{\tiny{[P1$\rangle$GM]}}
                \texttt{PYTANIE: Czy w twojej siatce w pierwszym i drugim wierszu od góry X znajduje się tylko w trzeciej kolumnie, czy może także w innych kolumnach?} \\
            }
        }
    }
    & & \\ \\

    \theutterance \stepcounter{utterance}  
    & & & \multicolumn{2}{p{0.3\linewidth}}{
        \cellcolor[rgb]{0.9,0.9,0.9}{
            \makecell[{{p{\linewidth}}}]{
                \texttt{\tiny{[GM$|$GM]}}
                \texttt{continue} \\
            }
        }
    }
    & & \\ \\

    \theutterance \stepcounter{utterance}  
    & & \multicolumn{4}{p{0.6\linewidth}}{
        \cellcolor[rgb]{0.9,0.9,0.9}{
            \makecell[{{p{\linewidth}}}]{
                \texttt{\tiny{[GM$\rangle$P2]}}
                \texttt{ODPOWIEDŹ: Tak, w piątym (dolnym) wierszu X znajduje się we wszystkich pięciu kolumnach.} \\
\texttt{PYTANIE: Czy w twojej siatce w pierwszym i drugim wierszu od góry X znajduje się tylko w trzeciej kolumnie, czy może także w innych kolumnach?Rozpocznij odpowiedź od "ODPOWIEDŹ:" i nie dodawaj nic więcej.} \\
            }
        }
    }
    & \\ \\

    \theutterance \stepcounter{utterance}  
    & & & \multicolumn{4}{p{0.6\linewidth}}{
        \cellcolor[rgb]{0.9,0.9,0.9}{
            \makecell[{{p{\linewidth}}}]{
                \texttt{\tiny{[GM$\langle$P2]}}
                \texttt{ODPOWIEDŹ: Tak, w pierwszym i drugim wierszu X znajduje się tylko w trzeciej kolumnie.} \\
            }
        }
    }
     \\ \\

    \theutterance \stepcounter{utterance}  
    & & & \multicolumn{2}{p{0.3\linewidth}}{
        \cellcolor[rgb]{0.9,0.9,0.9}{
            \makecell[{{p{\linewidth}}}]{
                \texttt{\tiny{[GM$|$GM]}}
                \texttt{continue} \\
            }
        }
    }
    & & \\ \\

    \theutterance \stepcounter{utterance}  
    & & \multicolumn{4}{p{0.6\linewidth}}{
        \cellcolor[rgb]{0.9,0.9,0.9}{
            \makecell[{{p{\linewidth}}}]{
                \texttt{\tiny{[GM$\rangle$P2]}}
                \texttt{Teraz zadaj pytanie, aby dowiedzieć się o nowych aspektach mojej siatki, które mogą różnić się od twojej siatki. Zacznij od "PYTANIE:" i nie dodawaj nic więcej.} \\
            }
        }
    }
    & \\ \\

    \theutterance \stepcounter{utterance}  
    & & & \multicolumn{4}{p{0.6\linewidth}}{
        \cellcolor[rgb]{0.9,0.9,0.9}{
            \makecell[{{p{\linewidth}}}]{
                \texttt{\tiny{[GM$\langle$P2]}}
                \texttt{PYTANIE: Czy w twojej siatce w czwartym wierszu od góry X znajduje się tylko w trzeciej kolumnie, czy także w czwartej?} \\
            }
        }
    }
     \\ \\

    \theutterance \stepcounter{utterance}  
    & & & \multicolumn{2}{p{0.3\linewidth}}{
        \cellcolor[rgb]{0.9,0.9,0.9}{
            \makecell[{{p{\linewidth}}}]{
                \texttt{\tiny{[GM$|$GM]}}
                \texttt{continue} \\
            }
        }
    }
    & & \\ \\

    \theutterance \stepcounter{utterance}  
    & & \multicolumn{4}{p{0.6\linewidth}}{
        \cellcolor[rgb]{0.9,0.9,0.9}{
            \makecell[{{p{\linewidth}}}]{
                \texttt{\tiny{[P1$\langle$GM]}}
                \texttt{ODPOWIEDŹ: Tak, w pierwszym i drugim wierszu X znajduje się tylko w trzeciej kolumnie.} \\
\texttt{PYTANIE: Czy w twojej siatce w czwartym wierszu od góry X znajduje się tylko w trzeciej kolumnie, czy także w czwartej?Rozpocznij odpowiedź od "ODPOWIEDŹ:" i nie dodawaj nic więcej.} \\
            }
        }
    }
    & \\ \\

    \theutterance \stepcounter{utterance}  
    & \multicolumn{4}{p{0.6\linewidth}}{
        \cellcolor[rgb]{0.9,0.9,0.9}{
            \makecell[{{p{\linewidth}}}]{
                \texttt{\tiny{[P1$\rangle$GM]}}
                \texttt{ODPOWIEDŹ: W mojej siatce w czwartym wierszu od góry X znajduje się tylko w trzeciej kolumnie.} \\
            }
        }
    }
    & & \\ \\

    \theutterance \stepcounter{utterance}  
    & & & \multicolumn{2}{p{0.3\linewidth}}{
        \cellcolor[rgb]{0.9,0.9,0.9}{
            \makecell[{{p{\linewidth}}}]{
                \texttt{\tiny{[GM$|$GM]}}
                \texttt{continue} \\
            }
        }
    }
    & & \\ \\

    \theutterance \stepcounter{utterance}  
    & & \multicolumn{4}{p{0.6\linewidth}}{
        \cellcolor[rgb]{0.9,0.9,0.9}{
            \makecell[{{p{\linewidth}}}]{
                \texttt{\tiny{[P1$\langle$GM]}}
                \texttt{Teraz zadaj pytanie, aby dowiedzieć się o nowych aspektach mojej siatki, które mogą różnić się od twojej siatki. Zacznij od "PYTANIE:" i nie dodawaj nic więcej.} \\
            }
        }
    }
    & \\ \\

    \theutterance \stepcounter{utterance}  
    & \multicolumn{4}{p{0.6\linewidth}}{
        \cellcolor[rgb]{0.9,0.9,0.9}{
            \makecell[{{p{\linewidth}}}]{
                \texttt{\tiny{[P1$\rangle$GM]}}
                \texttt{PYTANIE: Czy w twojej siatce w czwartym wierszu od góry X znajduje się także w czwartej kolumnie?} \\
            }
        }
    }
    & & \\ \\

    \theutterance \stepcounter{utterance}  
    & & & \multicolumn{2}{p{0.3\linewidth}}{
        \cellcolor[rgb]{0.9,0.9,0.9}{
            \makecell[{{p{\linewidth}}}]{
                \texttt{\tiny{[GM$|$GM]}}
                \texttt{continue} \\
            }
        }
    }
    & & \\ \\

    \theutterance \stepcounter{utterance}  
    & & \multicolumn{4}{p{0.6\linewidth}}{
        \cellcolor[rgb]{0.9,0.9,0.9}{
            \makecell[{{p{\linewidth}}}]{
                \texttt{\tiny{[GM$\rangle$P2]}}
                \texttt{ODPOWIEDŹ: W mojej siatce w czwartym wierszu od góry X znajduje się tylko w trzeciej kolumnie.} \\
\texttt{PYTANIE: Czy w twojej siatce w czwartym wierszu od góry X znajduje się także w czwartej kolumnie?Rozpocznij odpowiedź od "ODPOWIEDŹ:" i nie dodawaj nic więcej.} \\
            }
        }
    }
    & \\ \\

    \theutterance \stepcounter{utterance}  
    & & & \multicolumn{4}{p{0.6\linewidth}}{
        \cellcolor[rgb]{0.9,0.9,0.9}{
            \makecell[{{p{\linewidth}}}]{
                \texttt{\tiny{[GM$\langle$P2]}}
                \texttt{ODPOWIEDŹ: Tak, w mojej siatce w czwartym wierszu od góry X znajduje się zarówno w trzeciej, jak i w czwartej kolumnie.} \\
            }
        }
    }
     \\ \\

    \theutterance \stepcounter{utterance}  
    & & & \multicolumn{2}{p{0.3\linewidth}}{
        \cellcolor[rgb]{0.9,0.9,0.9}{
            \makecell[{{p{\linewidth}}}]{
                \texttt{\tiny{[GM$|$GM]}}
                \texttt{continue} \\
            }
        }
    }
    & & \\ \\

    \theutterance \stepcounter{utterance}  
    & & \multicolumn{4}{p{0.6\linewidth}}{
        \cellcolor[rgb]{0.9,0.9,0.9}{
            \makecell[{{p{\linewidth}}}]{
                \texttt{\tiny{[GM$\rangle$P2]}}
                \texttt{Teraz podejmij decyzję. Jak myślisz: czy Twoja siatka i druga opisana siatka są takie same? Napisz "DECYZJA: ta sama siatka" jeśli uważasz, że to ta sama siatka, albo "DECYZJA: różne siatki" jeśli uważasz, że to różne siatki. Nie dodawaj niczego więcej.} \\
            }
        }
    }
    & \\ \\

    \theutterance \stepcounter{utterance}  
    & & & \multicolumn{4}{p{0.6\linewidth}}{
        \cellcolor[rgb]{0.9,0.9,0.9}{
            \makecell[{{p{\linewidth}}}]{
                \texttt{\tiny{[GM$\langle$P2]}}
                \texttt{DECYZJA: różne siatki} \\
            }
        }
    }
     \\ \\

    \theutterance \stepcounter{utterance}  
    & & & \multicolumn{2}{p{0.3\linewidth}}{
        \cellcolor[rgb]{0.9,0.9,0.9}{
            \makecell[{{p{\linewidth}}}]{
                \texttt{\tiny{[GM$|$GM]}}
                \texttt{continue} \\
            }
        }
    }
    & & \\ \\

    \theutterance \stepcounter{utterance}  
    & & & \multicolumn{2}{p{0.3\linewidth}}{
        \cellcolor[rgb]{0.9,0.9,0.9}{
            \makecell[{{p{\linewidth}}}]{
                \texttt{\tiny{[GM$|$GM]}}
                \texttt{success} \\
            }
        }
    }
    & & \\ \\

    \theutterance \stepcounter{utterance}  
    & & \multicolumn{4}{p{0.6\linewidth}}{
        \cellcolor[rgb]{0.9,0.9,0.9}{
            \makecell[{{p{\linewidth}}}]{
                \texttt{\tiny{[P1$\langle$GM]}}
                \texttt{ODPOWIEDŹ: Tak, w mojej siatce w czwartym wierszu od góry X znajduje się zarówno w trzeciej, jak i w czwartej kolumnie.} \\
\texttt{Teraz podejmij decyzję. Jak myślisz: czy Twoja siatka i druga opisana siatka są takie same? Napisz "DECYZJA: ta sama siatka" jeśli uważasz, że to ta sama siatka, albo "DECYZJA: różne siatki" jeśli uważasz, że to różne siatki. Nie dodawaj niczego więcej.} \\
            }
        }
    }
    & \\ \\

    \theutterance \stepcounter{utterance}  
    & \multicolumn{4}{p{0.6\linewidth}}{
        \cellcolor[rgb]{0.9,0.9,0.9}{
            \makecell[{{p{\linewidth}}}]{
                \texttt{\tiny{[P1$\rangle$GM]}}
                \texttt{DECYZJA: różne siatki} \\
            }
        }
    }
    & & \\ \\

    \theutterance \stepcounter{utterance}  
    & & & \multicolumn{2}{p{0.3\linewidth}}{
        \cellcolor[rgb]{0.9,0.9,0.9}{
            \makecell[{{p{\linewidth}}}]{
                \texttt{\tiny{[GM$|$GM]}}
                \texttt{continue} \\
            }
        }
    }
    & & \\ \\

    \theutterance \stepcounter{utterance}  
    & & & \multicolumn{2}{p{0.3\linewidth}}{
        \cellcolor[rgb]{0.9,0.9,0.9}{
            \makecell[{{p{\linewidth}}}]{
                \texttt{\tiny{[GM$|$GM]}}
                \texttt{success} \\
            }
        }
    }
    & & \\ \\

\end{supertabular}
}

\end{document}
